\subsection{Matrice Densità}
Fino ad'ora abbiamo sempre supposto che nei sistemi presi in esame si partisse da degli stati ben precisi, che fossero autostati o combinazioni; ad esempio con Stern e Gerlach si aveva che il fascio in ingresso era già nello stato $\Ket {\psi_0} = \sqrthalf \l \Ket \uparrow + \Ket \downarrow \r$. Tuttavia  la situazione più generale deve rappresentare anche stati che non sono normalmente descrivibili come elementi degli spazi usati fin'ora (infatti, se si facesse passare $\Ket {\psi_0}$ come appena definito da $\SG{x}$ risulterebbe che il fascio è già polarizzato!). Per descrivere queste "miscele" di stati risulta comodo introdurre il concetto di matrice densità, così da tenere conto di sovrapposizioni più complesse di quelle solite.
\tab

La matrice densità è definita come segue. Si parte con una successione di 'pesi':
\[
    \omega_i \geq 0\ , \ \ \ \sum_{i=1}^{N} \omega_i = 1
\]
e con un operatore con la sua base:
\[
    \hat A \ , \ \ \ \{ \Ket {a_i} \} \ \ \implies \ \ \hat A \Ket {a_i} = a_i \Ket {a_i}  
\]
Si definisce la media statistica di $\hat A$ in modo diverso dal solito:
\[
    \overline A = \sum_{i=1}^{N} \omega_i \Bra {a_i} \hat A \Ket {a_i}
\]
Ad ora le $\omega_i$ sono i rapporti con cui si "miscelano" gli stati. Alcuni chiamano questa miscela "combinazione classica". Non si tratta di sovrapposizione in senso quantistico. Questo risultato è espandibile in termini di una nuova base (ortonormale) a piacere $\{ \Ket b \}$:
\[
    \overline A = \sum_{i=1}^{N} \omega_i \Bra {a_i} \sum_{b} \ketbra{b}{b} \hat A \sum_{b'} \ketbra{b'}{b'} \Ket {a_i} = \]\[
    \sum_{b,b'} \Bra {b'} \left( \sum_{i=1}^{N} \omega_i \ketbra{a_i}{a_i} \right) \Ket {b} \Bra {b} \hat A \Ket {b'} = \sum_{b'}  \Bra {b'} \hat \rho \hat A \Ket {b'}
\]
Evidentemente la quantità finale è la traccia di $\hat \rho A$, ed è indipendente dalla scelta di base. Qeusto operatore $\hat \rho$ è detto \textbf{matrice densità}.\index{Matrice Densità} Tra le proprietà della matrice associata all'operatore si ha che è hermitiana, ha traccia uguale a uno ed è definita positiva. Evidentemente la traccia è anche la somma dei suoi autovalori (tutti reali).

Questa matrice densità è evidentemente diagonale quando espressa nella base di $\A$.

Se per esempio si avesse che $\omega_i = \delta_{0i}$ che il sistema è in un autostato di $\A$, quindi si ricade nei casi analizzati precedentemente.

\subsubsection*{Spazio Puro}
Un operatore $\hr$ si dice \textbf{Stato Puro}\index{Stati puri} se esiste uno stato $\Ket \psi$ generico tale che la matrice densità è scrivibile come $\hat \rho = \ketbra \psi \psi$. Si ha anche che $\hr^2 = \hr$, al chè si può mostrare che in questo caso gli autovalori possono essere solo 0 e 1.\\
Lo spazio puro è lo spazio degli stati puri.

\subsubsection*{Esempio 1: $\Ket {\psi_0} = \Up$}
Risulta:
\[
    \hr = \ketbrabra{\psi_0} = \Up \Pu = \begin{pmatrix}
        1 & 0 \\
        0 & 0 \\
    \end{pmatrix}
\]
che evidentemente rispetta tutti i requisiti.

\subsubsection*{Esempio 2: $\Ket {\psi_0} = \Ket {x_+}$}
Risulta:
\[
    \hat \rho = \ketbrabra{\psi_0} = \half \ketbra {x+} {x+} + \half \ketbra {x-} {x-} = \half \begin{pmatrix}
        1 & 1 \\
        1 & 1 \\
    \end{pmatrix}
\]
che evidentemente rispetta tutti i requisiti.

\subsubsection*{Esempio 3: Fascio non polarizzato}
Se rivediamo l'esperimento di Stern e Gerlach alla luce delle interpretazioni introdotte in questo capitolo, si ha che un fascio di elettroni non polarizzato risulta (statisticamente) in metà degli elettroni risultanti in $\Up$ e metà $\Dn$. Questa combinazione viene trascritta così come viene:
\[
    \hat \rho = \half \Up \Pu + \half \Dn \Nd = \half \begin{pmatrix}
        1 & 0 \\
        0 & 1 \\
    \end{pmatrix}
\]
che \textit{non} rispetta le richieste di stato puto (ma rispetta $\Tr(\hr) = 1$).

\subsubsection*{Esempio 4: Un altro fascio non polarizzato}
Stavolta scriviamo il fascio non polarizzato come composizione di $\Ket {x_+}$ e $\Ket {x_-}$ invece che $\Up$ e $\Pu$. Risulta:
\[
    \hat \rho = \half \ketbrabra{x_+} + \half \ketbrabra{x_-} = \half \begin{pmatrix}
        1 & 0 \\
        0 & 1 \\
    \end{pmatrix}
\]
che è uguale al risultato dell'Esempio 3, nonostante la costruzione è stata fatta in modo diverso! Questo suggerisce che non vi è corrispondenza biunivoca tra matrice densità e "stato" del sistema inteso come condizione classica in cui si trova il sistema.

\subsection{State purity}\index{State puriry}
Si può notare che negli esempi 3 e 4 si ha $\Tr(\hr^2) = 1/2 \leq 1$. Si può verificare che in generale è vero per ogni stato non puro, e l'equivalenza è vera solo per gli stati puri. Viene quindi naturale chiamare "State puriry" la quantità $\Tr(\hr^2)$, che ha massimo solo quanto lo sato è puro.
\tab

Se ci si restringe ai sistemi con spin $1/2$ allora si può fare un'analisi sistematica della matrice densità. Risulta, nel caso più generico possibile:
\[
    \hr = \begin{pmatrix}
        a_0 + ib_0 & a_1 + ib_1 \\
        a_3 + ib_3 & a_4 + ib_4 \\
    \end{pmatrix}
\]
Ora, ricordando del fatto che $\hr$ è hermitiano, si passa da 8 grai a solo 4. Se poi si ricorda che deve avere traccia unitaria risulta che:
\[
    \hr = \begin{pmatrix}
        a_0 & a_1 + ib_1 \\
        a_1 - ib_1 & 1 - a_0 \\
    \end{pmatrix}
\]
dipende solo da tre parametri reali. È possibile ricondurre i valori di questi parametri ai valori medi di $\SG{x}, \SG{y}$ e $\SG{z}$.
\tab

Ora ci si può porre la domanda: è possibile risalire a posteriori (tramite misurazioni di osservabili) al valore della matrice densità? In particolare, è possibile distringuere gli stati puri dagli stati mixed? Che tipo di misurazioni bisogna fare?

Partiamo dal calcolare il valor medio di $\SG{x}$ ed $\SG{z}$ nei casi riportati negli esempi 2 e 4. Il primo caso ricade nella situazione solita in cui non serve interpellare la matrice densità:
\[
    \overline{\SG{x}} = \braketin{x_+}{\SG{x}}{x_+} = \frac{\h}{2}\ , \ \ \ \overline{\SG{z}} = \braketin{x_+}{\SG{z}}{x_+} = 0
\]
Nel secondo caso invece risulta:
\[
    \overline{\SG{x}} = \Tr(\hr\SG{x}) = 0\ , \ \ \ \overline{\SG{z}} = \Tr(\hr\SG{z}) = 0
\]

In sostanza si ha che misurando $\SG{x}$ i due casi sono distinguibili, invece misurando $\SG{z}$ no! In sostanza non è possibile risalire alla matrice densità misurando il sistema con un solo osservabile.

\subsection{Misure parziali}\index{Misure parziali}
È importante analizzare alcuni particolari casi di misure, in quanto risulta che esistono misurazioni che non conservano la state puriry: un esempio sono le misure parziali.

Sia un sistema descritto dal prodotto tensoriale di due spazi di Hilbert $\Hilb_1$ e $\Hilb_2$. Uno stato generico in una base generica sarà del tipo $\Kp[0] = \sum_{n,m}C_{nm}\Ket {n,m}$ con $\Ket {n,m} = \Ket n \otimes \Ket m$. Per comodità si userà la lettera $n$ per indicare stati appartenenti al primo spazio e $m$ per stati nel secondo.\\
Allora la matrice densità sarà del tipo:
\[
    \hr = \ketbrabra{\psi_0} = \sum_{n,m}\sum_{n',m'} C_{nm}C^*_{n'm'}\ketbra{n,m}{n',m'}
\]
Sia ora $\A_1$ un operatore che agisce \textit{solo} sullo spazio $\Hilb_1$. È possibile promuoverlo a operatore in $\Hilb_1 \times \Hilb_2$ semplicemente definendo $\A = \A_1 \otimes \Id$, ovvero facendo agire $\A_0$ su $\Hilb_1$ e l'identità su $\Hilb_2$.\\
Se ora si calcola la media di $\A$:
\[
    \overline{\A} = \Tr(\hr\A) = \sum_n \sum_{m} C_{nm}C^*_{nm} \braketin{n,m}{\A}{n,m} = \sum_n \sum_{m} C_{nm}C^*_{nm} \braketin{n}{\A}{n}
\]
Ora, per comprendere meglio questo risulato è opportuno introdurre un nuovo operatore, chiamato Matrice densità ridotta\index{Matrice densità ridotta}. È definita come segue:
\[
    \hr_1 = \Tr_2(\hr) = \sum_{n,m}\sum_{n'} C_{nm}C^*_{n'm} \ketbra{n}{n'}
\]
dove $\Tr_2$ indica l'operazione traccia fatta solo sullo spazio $\Hilb_2$.
È facile verificare che rispetta tutte le ipotesi necessarie per essere una matrice densità, e per di più:
\[
    \Tr(\hr\A) = \Tr_1(\hr_1\A_1) =  \sum_{n,m} C_{nm}C^*_{nm} \ketbra{n}{n'}
\]
È opportuno verificare con un esempio il fatto che la matrice densità ridotta potrebbe essere non pura:
\subsubsection*{Esempio: elettrone e protone}
Si parte dallo stato puro $\Ket \psi = \sqrthalf \l \Up _p \otimes \Dn _e + \Dn _p \otimes \Up _e \r = \sqrthalf\l \Ket {\u,\d} + \Ket {\d,\u} \r$. Se si vuole calcolare ma media del momento angolare dell'elettrone in una direzione generica si calcola:
\[
    \overline{\S_{\vec n}} = \Tr\l \hr_e \hat\sigma_{\vec n} \r
\]
al ché:
\[
    \hr_e = \Tr_p (\hr) = \Tr_p \l \ketbrabra{\psi} \r = \Tr_p \left[\half \l \Ket {\u,\d} + \Ket {\d,\u} \r\l \Bra {\u,\d} + \Bra {\d,\u} \r \right] \sim \half \begin{pmatrix}
        1 & 0 \\ 0 & 1
    \end{pmatrix}
\]
Che non è puro!

\tab
È possibile quindi definire un criterio sufficiente per far sì che una misura parziale (dopotutto tutte le misure sono parziali, visto che ogni sistema è contenuto in un sistema più grande!) non sia soggetta a questo effetto? La risposta è \textbf{sì}: si può dimostrare che lo stato puro è conservato \textit{se e solo se} non vi è entanglement tra i sottosistemi in analisi, ovvero quando lo stato del sistema è scrivibile nella forma:
\[
    \Ket \psi = \bigotimes_n \Ket {\psi_n}
\]